% tk-04：五类核心转化示意（动→静 / 形→数 / 数→形 / 未知→已知 / 分散→集中）
\documentclass[tikz,border=4pt]{standalone}
\usepackage{ctex}
\usepackage{amsmath}
\usetikzlibrary{calc, arrows.meta, decorations.markings}
\begin{document}
\begin{tikzpicture}[scale=0.8, thick,
  myarr/.style={-{Latex[length=2.2mm]}, thick},
  tick/.style={postaction={decorate, decoration={markings,
    mark=at position 0.5 with {\draw (-1.5pt,-3pt) -- (1.5pt,3pt);}}}}]

  % 1. 动 → 静
  \begin{scope}[xshift=0cm, yshift=0cm]
    \draw (0,0) -- (3.4,0);
    \node[below left]  at (0,0)   {\footnotesize $A$};
    \node[below right] at (3.4,0) {\footnotesize $B$};
    % 动点 P 的三个位置：浅灰 + 实点
    \fill[gray!50] (0.8,0) circle (1.4pt);
    \fill[gray!50] (1.6,0) circle (1.4pt);
    \fill[red] (2.4,0) circle (1.8pt);
    \node[above] at (0.8,0) {\footnotesize $P_1$};
    \node[above] at (1.6,0) {\footnotesize $P_2$};
    \node[above] at (2.4,0) {\footnotesize $P(t)$};
    \draw[myarr, gray] (0.8,-0.05) -- (2.4,-0.05);
    \node[below] at (1.7,-1.0) {\small 动 $\to$ 静（设参数）};
  \end{scope}

  % 2. 形 → 数
  \begin{scope}[xshift=6cm, yshift=0cm]
    \draw[->] (-0.3,0) -- (3.2,0) node[right] {\footnotesize $x$};
    \draw[->] (0,-0.3) -- (0,2.4) node[above] {\footnotesize $y$};
    \coordinate (A) at (0.6,0.4);
    \coordinate (B) at (2.6,1.8);
    \fill (A) circle (1.5pt);
    \fill (B) circle (1.5pt);
    \draw (A) -- (B);
    \node[below left] at (A) {\footnotesize $A(x_1,y_1)$};
    \node[above right] at (B) {\footnotesize $B(x_2,y_2)$};
    \node[below] at (1.4,-1.0) {\small 形 $\to$ 数（建坐标）};
  \end{scope}

  % 3. 数 → 形
  \begin{scope}[xshift=12cm, yshift=0cm]
    \node[align=center] at (1.6,1.8) {\footnotesize $\sqrt{x^2+1}$\\[-1pt]\footnotesize $+\sqrt{(x-3)^2+4}$};
    \draw[myarr] (1.6,1.1) -- (1.6,0.6);
    \draw[->] (-0.2,0) -- (3.4,0);
    \coordinate (A) at (0,0.5);
    \coordinate (B) at (2.4,1.0);
    \coordinate (P) at (1.0,0);
    \fill (A) circle (1.2pt);
    \fill (B) circle (1.2pt);
    \fill[red] (P) circle (1.3pt);
    \draw[red] (A) -- (P) -- (B);
    \node[left] at (A) {\footnotesize $A$};
    \node[above right] at (B) {\footnotesize $B$};
    \node[below] at (P) {\footnotesize $P$};
    \node[below] at (1.5,-1.0) {\small 数 $\to$ 形（看距离）};
  \end{scope}

  % 4. 未知 → 已知（搭桥）
  \begin{scope}[xshift=0cm, yshift=-5cm]
    \node[draw, circle, minimum size=0.9cm] (U) at (0,0.5)   {\footnotesize 未知};
    \node[draw, circle, minimum size=0.9cm] (M) at (1.7,0.5) {\footnotesize 桥};
    \node[draw, circle, minimum size=0.9cm] (K) at (3.4,0.5) {\footnotesize 已知};
    \draw[myarr, red] (U) -- (M);
    \draw[myarr, red] (M) -- (K);
    \node[below] at (1.7,-1.0) {\small 未知 $\to$ 已知（搭桥）};
  \end{scope}

  % 5. 分散 → 集中
  \begin{scope}[xshift=6cm, yshift=-5cm]
    % 左：分散
    \draw[tick] (-0.3,1.5) -- (0.7,1.5);
    \draw[tick] (1.6,0.2) -- (2.3,0.7);
    \node at (0.2,1.85) {\footnotesize $a$};
    \node at (2.1,0.95) {\footnotesize $b$};
    % 箭头
    \draw[myarr] (2.8,1.0) -- (3.6,1.0);
    % 右：集中（首尾相接的三角形）
    \draw[tick] (3.9,0.2) -- (4.9,0.2);
    \draw[tick] (4.9,0.2) -- (5.4,1.0);
    \draw[dashed] (3.9,0.2) -- (5.4,1.0);
    \node[below] at (4.4,0.15) {\footnotesize $a$};
    \node[right] at (5.15,0.6) {\footnotesize $b$};
    \node[below] at (2.5,-1.0) {\small 分散 $\to$ 集中（平移/旋转）};
  \end{scope}

\end{tikzpicture}
\end{document}
